\section{Conclusion}
A generic IR may be structurally valid while violating the semantic contracts of a selected compiler composition. Obligation-guided reverification turns declared invalidations into verifier-owned duties that must be discharged by a due boundary, independent of later consumer demand. Under explicit completeness and soundness assumptions, a successful boundary transition cannot leave a registered due obligation pending.

Across a frozen System W boundary corpus, a 30-program source-to-result study, and a public-SDK TensorRules package, selective obligation discharge matches always verification on evaluated classifications and rejects faults missed by structural and invalidation-only policies. A separate lazy-demand witness isolates the key trigger difference, while eight mechanism ablations provide concrete necessity witnesses. The evidence supports relative boundary safety, detection, and localization for the evaluated contracts. It does not establish general verifier completeness, natural fault prevalence, or whole-compilation speedup; externally authored faults and pinned-machine performance remain necessary before making those claims.
